% TOP_PROBLEM: {"id": 30004208, "title": "The Dixmier Unitarisability Problem", "category_id": 4, "category": {"id": 4, "name": "algebra", "display_name": "Algebra", "description": "Group theory, ring theory, field theory, and algebraic structures.", "slug": "algebra", "order_index": 4, "created_at": "2026-07-31T15:26:25.671Z"}, "status": "open"}
% FIRST_POSTED: 2026-09-25
% ENTRY_KIND: statement_only
% !TeX program = lualatex
% Definition and sources only; this file is not an LLM solution attempt.
\documentclass[11pt]{article}
\usepackage[margin=25mm]{geometry}
\usepackage{fontspec}
\setmainfont{Latin Modern Roman}
\usepackage{amsmath,amssymb,mathtools}
\usepackage{unicode-math}
\setmathfont{Latin Modern Math}
\usepackage{xurl,hyperref}
\hypersetup{colorlinks=true,urlcolor=blue}
\setcounter{secnumdepth}{0}
\setlength{\parindent}{0pt}
\setlength{\parskip}{5pt}
\setlength{\emergencystretch}{3em}
\begin{document}
\section*{\texorpdfstring{The Dixmier Unitarisability Problem}{The Dixmier Unitarisability Problem}}\label{30004208}
\subsection{Definitions and mathematical statement}
A discrete group $G$ is unitarisable if every homomorphism $\pi:G\to\operatorname{GL}(\mathcal H)$ on every complex Hilbert space, with $\sup_g\|\pi(g)\|<\infty$, has a bounded invertible operator $S$ such that
\[
 (S\pi(g)S^{-1})^*(S\pi(g)S^{-1})=I\qquad(g\in G).
\]
Amenability means existence of a positive normalized left-invariant mean on $\ell^\infty(G)$. The conjectured implication is that unitarisability forces amenability; equivalently, every nonamenable discrete group has at least one uniformly bounded representation not similar to a unitary representation. No countability assumption on $G$, separability assumption on $\mathcal H$, or prescribed similarity bound on $S$ is included. Boundedness of each individual $\pi(g)$ is weaker than the required uniform bound over the group.
\par\smallskip\textit{Formulation and concept references:} \href{https://arxiv.org/pdf/0902.4585v1}{[S1]}.

\subsection{Short English statement}
If every uniformly bounded Hilbert-space representation of a discrete group can be made unitary by a change of coordinates, must the group be amenable?
\subsection{Sources}
\begin{itemize}
\item[S1]Nicolas Monod and Narutaka Ozawa, The Dixmier problem, lamplighters and Burnside groups, arXiv0902.4585v1 (26Feb2009).
\url{https://arxiv.org/pdf/0902.4585v1}.
\textit{Formulation source.}
\end{itemize}
\end{document}
